% ============================================================
%  Chapter 1: Introduction and Objective
%  (descriptive title retained per author preference; the FAPS
%  template names this chapter "Introduction" -- either is fine.)
% ============================================================
\chapter[Introduction]{Introduction to Leading Indicators for Operational Safety Assurance}
\label{ch:introduction}

% Descriptive chapter titles follow supervisor feedback (P. Ziegler): the ToC
% shows the full title, the optional argument keeps the running head short.

This chapter sets out why the operational safety assurance of an automated driving
system depends on indicators that anticipate a loss event rather than record one, and
why the construction of such indicators is not yet a solved problem. It states the
objective of the thesis, the three research questions that structure the work, the
scenario on which the work is validated, and the organisation of the chapters that
follow.

\section{Motivation}
\label{sec:intro_motivation}

An automated driving system reaches public roads only once a safety case has argued
that it is acceptably safe for a stated operating domain, built around the kind of
structured, evidence-supported argument introduced in Chapter~\ref{ch:fundamentals}.
That argument is constructed and reviewed before deployment, yet its subject keeps
producing evidence about itself only once it is operating among road users whose
behaviour was never fully enumerable at design time. Some of the doubts an engineer
raises while building the argument cannot be closed through development-time testing
at all, because the evidence that would close them is precisely the evidence that
operation itself produces. A safety case for such a system is therefore not a document
that is completed once at approval. It remains open on specific, identifiable points,
and those points stay open until operational evidence closes them.

The scale of the difficulty is not small. Demonstrating, from observed outcomes alone,
that an automated driving system achieves a fatality rate meaningfully better than the
human baseline would require, by established estimates, on the order of hundreds of
millions to billions of kilometres of driving \cite{kalra2016driving}, a volume no
fleet reaches within a commercially relevant timeframe. Waiting for lagging evidence,
meaning evidence recorded only once a hazardous outcome has already occurred, is
consequently not a workable basis for closing the open points in an operational safety
case. Both ANSI/UL~4600 and PAS~1881 address this directly: a deployed system is
required to be monitored through indicators that anticipate a hazardous outcome rather
than merely record it, and reliance on lagging indicators alone is explicitly ruled out
\cite{ul4600, pas1881}. Koopman and Wagner make a closely related argument from the
deployment decision itself, holding that a defensible decision to deploy rests on a
positive trust balance the developer commits to sustaining after release, rather than
on a body of pre-deployment evidence considered sufficient on its own
\cite{koopman2020positive}.

What neither standard supplies, and what the review in Chapter~\ref{ch:sota} shows the
wider literature does not supply either, is a procedure for constructing such an
indicator in the first place. Existing work is well developed at locating where in a
safety argument an indicator is needed and at specifying how a system should respond
once one is violated, but the step in between, deriving the indicator's signal, source
and threshold from a specific, unresolved doubt in the argument, is left to engineering
judgement exercised case by case. This thesis develops a systematic procedure for that
step, the Causal Chain Derivation Method (CCDM) presented in Chapter~\ref{ch:methodology},
which takes an unresolved doubt surviving the eliminative elaboration of a Goal
Structuring Notation claim and derives from it a leading indicator whose signal, source
and threshold are all traceable back to the claim it addresses. The method is developed
generically and then validated empirically against the criteria set out in the same
chapter.

\section{Problem Statement}
\label{sec:intro_problem}

The gap that this thesis addresses is therefore narrow and specific. It is not the
absence of a requirement to monitor, which the standards already impose, nor the
absence of a place in the argument at which an indicator belongs, which the
assurance-case literature already supplies. It is the absence of a reproducible
procedure that turns a named, unresolved doubt into a named, calibrated quantity. In
the absence of such a procedure the derivation is performed anew by each engineer for
each claim, which makes the resulting indicator set neither reviewable nor
transferable, and leaves the connection between an indicator and the claim it is
meant to support asserted rather than demonstrated.

The direction of current vehicle regulation reinforces why this is worth solving
systematically rather than case by case. Germany's ordinance governing the operation of
vehicles with autonomous driving functions already imposes operational monitoring
obligations on deployed systems, and the assessment method under development at the
United Nations Economic Commission for Europe points toward comparable in-service
monitoring and reporting expectations at the international level \cite{afgbv2022,
natm2021}. A method that derives leading indicators traceably from a safety case,
rather than one that depends on an individual engineer's judgement repeated anew for
every claim, is accordingly not only a methodological convenience but a step toward
meeting an obligation that regulation is moving to require outright.

\section{Objective and Research Questions}
\label{sec:intro_objective}

The objective of this thesis is to develop this derivation procedure and to establish
whether the indicators it produces are usable for the operational safety assurance of
an automated driving system. Three research questions structure the work. The first
asks how a leading safety performance indicator, comprising a signal, a source, and a
threshold, can be derived systematically and traceably from an unresolved doubt in a
Goal Structuring Notation safety case. The second asks whether the indicators so
derived detect a developing violation early enough to permit a response, rather than
merely confirming that a violation has occurred. The third asks whether a single
derivation method transfers across conflicts of different character, namely
longitudinal, crossing, and following conflicts, within one manoeuvre. The first
question is addressed by the method of Chapter~\ref{ch:methodology}, and the second and
third by its empirical evaluation in Chapter~\ref{ch:results}.

\section{Scope of the Empirical Validation}
\label{sec:intro_scope}

The empirical validation is carried out on an unprotected left turn at an urban
intersection, in which the ego vehicle must negotiate oncoming traffic and a crossing
pedestrian without a protected signal phase, while separately maintaining a safe
separation from traffic following behind it. The left turn across the path of oncoming traffic is treated in the literature as
one of the most demanding interactions an automated vehicle meets at an
intersection, as reviewed in Section~\ref{sec:sota_conflicts}. The scenario is a demanding one for a single
derivation method to be tested against, precisely because it folds together three
conflicts of different character within one manoeuvre: a longitudinal conflict with
oncoming traffic, a crossing conflict with a pedestrian, and a longitudinal following
conflict to the rear. Whether indicators derived by the same method transfer across
these three geometries, or whether they do not, is itself informative about the
assumptions the method makes, and Chapter~\ref{ch:results} reports this comparison
directly.

\section{Structure of the Thesis}
\label{sec:intro_structure}

The remainder of this thesis is structured as follows. Chapter~\ref{ch:fundamentals}
establishes the terminology of safety assurance, safety argumentation, and safety
performance indicators, together with the surrogate safety metrics and the simulation
environment used in the later chapters. Chapter~\ref{ch:sota} reviews the current state
of research, covering operational safety assurance, accident causation and
investigation, the transition from reactive to proactive monitoring, and the derivation
of leading indicators from lagging ones across safety domains, and it closes with the
research gap that this work addresses. Chapter~\ref{ch:methodology} presents the Causal
Chain Derivation Method and its worked derivations for the three conflicts of the use
case. Chapter~\ref{ch:experiment} describes the simulation-based evaluation built on
the CARLA simulator, including the scenario, the sampled parameters, the collision
definition and the evaluation measures. Chapter~\ref{ch:results} reports the findings
of that evaluation on the unprotected left-turn scenario. Chapter~\ref{ch:discussion} interprets the
results with respect to the objective and the state of the art and states the
limitations of the work, and Chapter~\ref{ch:conclusion} summarises the contributions
and outlines directions for future work.